% ==STYLE-BEGIN==
\documentclass[UTF8,11pt]{ctexart}
\usepackage{amsmath,amssymb,amsthm}
\usepackage[a4paper,margin=2.4cm,top=2.3cm,bottom=2.3cm]{geometry}
\usepackage{fancyhdr}
\usepackage{enumitem}
\usepackage{hyperref}
\hypersetup{hidelinks}

% ---- 版式 ----
\linespread{1.15}
\setlength{\parskip}{0.35em}
\setlength{\parindent}{1.6em}
\setlist{nosep,leftmargin=2em}

% ---- 页眉 ----
\pagestyle{fancy}
\fancyhf{}
\fancyhead[L]{\small\coursename}
\fancyhead[R]{\small\thepage}
\renewcommand{\headrulewidth}{0.4pt}

% ---- 定理环境：中文用正体，标题加粗，按小节共用一条流水号 ----
\newtheoremstyle{cnthm}{0.7em}{0.7em}{\normalfont}{}{\bfseries}{}{0.5em}{}
\theoremstyle{cnthm}
\newtheorem{thm}{定理}[section]
\newtheorem{lem}[thm]{引理}
\newtheorem{prop}[thm]{命题}
\newtheorem{cor}[thm]{推论}
\newtheorem{defn}[thm]{定义}
\newtheorem{rem}[thm]{注}
\newtheorem{exa}[thm]{例}
\newtheorem{conj}[thm]{猜想}
\renewcommand{\proofname}{\bfseries 证明}
\providecommand{\coursename}{}
% ==STYLE-END==
\renewcommand{\coursename}{数学分析\,·\,1012}
\title{数学分析 1012}
\date{}
\begin{document}
\maketitle

\section{Recall}

函数的极限分为两部分,函数值的趋势与自变量的趋势
数列的极限一般考虑趋向正/负无穷
而函数的极限考虑趋向正/负无穷,趋向一点,从左/右趋向一点,即
$$
x \to +\infty \quad x \to -\infty \quad x \to x_0 \quad x \to x_0^+ \quad x \to x_0^-
$$
它们会对应函数极限的定义,即
\[
\begin{aligned}
&\forall\varepsilon>0,\ \exists N>0,\ \forall x>N,\quad
|f(x)-A|<\varepsilon &&(x\to+\infty),\\
&\forall\varepsilon>0,\ \exists N>0,\ \forall x<-N,\quad
|f(x)-A|<\varepsilon &&(x\to-\infty),\\
&\forall\varepsilon>0,\ \exists\delta>0,\ \forall x,\quad
0<|x-x_0|<\delta\Longrightarrow |f(x)-A|<\varepsilon
&&(x\to x_0),\\
&\forall\varepsilon>0,\ \exists\delta>0,\ \forall x,\quad
0<x-x_0<\delta\Longrightarrow |f(x)-A|<\varepsilon
&&(x\to x_0^+),\\
&\forall\varepsilon>0,\ \exists\delta>0,\ \forall x,\quad
0<x_0-x<\delta\Longrightarrow |f(x)-A|<\varepsilon
&&(x\to x_0^-).
\end{aligned}
\]
\section{函数极限的性质}

1. 极限存在即唯一
$$
\text{if} \lim\limits_{x \to x_0}f(x) = A \  \text{\&} \lim\limits_{x \to x_0}f(x) = B \to A=B
$$
证明
将极限写成定义形式有
$$\begin{aligned}
\lim\limits_{x \to x_0}f(x) = A \to \forall \varepsilon > 0 , \exists \delta > 0 , \forall x \ s.t. \ 0<|x-x_0|< \delta , \ s.t. \ |f(x)-A|< \frac{\varepsilon}{2} \\
\lim\limits_{x \to x_0}f(x) = B \to \forall \varepsilon > 0 , \exists \delta > 0 , \forall x \ s.t. \ 0<|x-x_0|< \delta , \ s.t. \ |f(x)-B|< \frac{\varepsilon}{2}
\end{aligned}
$$
之前我们已经证明过,$\varepsilon$前有任意定系数不影响极限,我们这里取前系数为$\frac{1}{2}$
考虑$|A-B|$有
$$
|A-B| \leq |A-f(x)|+|f(x)-B| < \frac{\varepsilon}{2}+\frac{\varepsilon}{2} = \varepsilon
$$
即A=B
2. 局部有界性
$$
\text{if} \lim\limits_{x \to x_0}f(x) = A , \exists M > 0 , \exists \delta > 0 \ s.t. \ \forall x (0<|x-x_0|< \delta) , |f(x)| \leq M
$$
回顾数列的有界性,我们知道"收敛数列必有界". 在数列中,我们对于$n>N$的无限多项有界,而前面是有限多项,所以整个数列必然有界
而在函数中,我们仍然可以得到$x_0$附近的函数是有界的,但是在这之外函数有无穷多个点,我们只能控制在$x_0$附近的点,故称"局部有界性"
证明
依旧利用极限的定义
$$
\lim\limits_{x \to x_0}f(x) = A \to \forall \varepsilon > 0 , \exists \delta > 0 , \forall x \ s.t. \ 0<|x-x_0|< \delta , \ s.t. \ |f(x)-A|< \varepsilon
$$
此时,我们取一个特定的$\varepsilon = 1$有
$$\begin{aligned}
\exists \delta_{1} > 0 , \forall x (0<|x-x_0|<\delta) , |f(x)-A|<1 \\
\Rightarrow f(x)| \leq max \{ |A|+1 , |A|-1 \}
\end{aligned}
$$
3. 保号性
$$
\text{if} \lim\limits_{x \to x_0}f(x) = A > 0 , \exists \delta > 0 , \forall x (0<|x-x_0|<\delta)) , f(x) > \frac{1}{2}A > 0
$$
取$\varepsilon = \frac{1}{2} A$易证
4. 四则运算
$$
\lim\limits_{x\to x_0}f(x)=A,\qquad
\lim\limits_{x\to x_0}g(x)=B
$$
则有
$$
\lim\limits_{x\to x_0}\bigl(f(x)\pm g(x)\bigr)=A\pm B,
\qquad
\lim\limits_{x\to x_0}\bigl(f(x)g(x)\bigr)=AB
$$
若 $B\neq0$，则
$$
\lim\limits_{x\to x_0}\frac{f(x)}{g(x)}
=\frac{A}{B}
$$
\begin{quote}
参与四则运算的极限要是同一个极限过程
\end{quote}


\section{海涅原理}

定理:
$$
\lim\limits_{x \to x_0}f(x) = A
\Leftrightarrow
\forall x_n \to x_0 (x_n \neq x_0)
\text{数列}\{ f(x_n) \}\text{有}
\lim\limits_{n \to +\infty}f(x_n) = A
$$
与归一原理类似,海涅原理完成了整体与局部,连续与离散的转换
\begin{quote}
$\{ f(x_n) \}$ 是一个数列 , 第n项的值是$f(x_n)$
$x_n \to x_0$ , $\{ x_n \}$本身也是数列
\end{quote}


证明
$\Rightarrow$ 由$\lim\limits_{x\to x_0}f(x)=A$可知
$$
\forall\varepsilon>0,\ \exists\delta>0,\ 
\forall x,\quad
0<|x-x_0|<\delta
\Longrightarrow
|f(x)-A|<\varepsilon.
$$
任取数列 $\{x_n\}$，$(x_n \to x_0)$ 寻找N , $\forall n , n > N , |f(x)-A| < \varepsilon$
要证明$\lim\limits_{n\to+\infty}f(x_n)=A$
只需证明
$$
\forall\varepsilon>0,\ \exists N,\ \forall n>N,
\quad
|f(x_n)-A|<\varepsilon.
$$
给定 $\varepsilon>0$，由函数极限的定义，存在 $\delta>0$，使得
$$
0<|x-x_0|<\delta
\Longrightarrow
|f(x)-A|<\varepsilon.
$$
又因为 $x_n\to x_0$，所以对于上述 $\delta>0$，存在 $N$，使得当 $n>N$ 时，
$$
|x_n-x_0|<\delta.
$$
同时，由于 $x_n\neq x_0$，有
$$
0<|x_n-x_0|<\delta.
$$
因此$|f(x_n)-A|<\varepsilon$
故
$$
\lim\limits_{n\to+\infty}f(x_n)=A
$$

$\Leftarrow$用反证法
假设
$$
\lim\limits_{x\to x_0}f(x)\neq A.
$$
则存在 $\varepsilon>0$，使得对于任意 $\delta>0$，都存在 $x$，满足
$$
0<|x-x_0|<\delta,
\qquad
|f(x)-A|\geq\varepsilon.
$$

特别地，取
$$
\delta=\frac{1}{n},
$$
则对于任意 $n\in\mathbb{N}$，存在 $x_n$，使得
$$
0<|x_n-x_0|<\frac{1}{n},
\qquad
|f(x_n)-A|\geq\varepsilon.
$$

由
$$
0<|x_n-x_0|<\frac{1}{n}
$$
可知
$$
\lim\limits_{n\to+\infty}x_n=x_0.
$$

但是
$$
|f(x_n)-A|\geq\varepsilon,
$$
所以
$$
\lim\limits_{n\to+\infty}f(x_n)\neq A.
$$

这与题设“对于任意满足 $x_n\neq x_0$ 且 $x_n\to x_0$ 的数列 $\{x_n\}$，都有 $f(x_n)\to A$矛盾。

因此
$$
\lim\limits_{x\to x_0}f(x)=A.
$$

e.g. 极限
$$
\lim\limits_{x \to 0} sin \frac{1}{x}
$$
不存在
分析
$sin\frac{1}{x_1} = sin\frac{1}{x_2} \Leftarrow \frac{1}{x_1} - \frac{1}{x_2} = 2\pi$
据此,我们可以得到
$$
|x_2-x_1| = 2\pi x_1x_2 < 2\pi x_2^2
$$
随着$x_2 \to 0$,$|x_2-x_1| \to 0$,震荡变得越来越激烈,直观上看没有极限
我们用海涅原理来证明,只需要构造出两条点列,当两个点列都有$n \to +\infty$时,$a_n$的趋势不一

证明
构造
$$
x_n = \frac{1}{2n\pi} \quad y_n = \frac{1}{2n\pi + \frac{\pi}{2}}
$$
我们有
$$\begin{aligned}
\sin{\frac{1}{x_n}} = \sin{2n\pi} = 0 \\
\sin(\frac{1}{y_n}) = \sin{2n\pi + \frac{\pi}{2}} = 1
\end{aligned}
$$
故极限不存在

\section{夹迫定理}

定理:
$$\begin{aligned}
\text{if} \ \exists \delta>0 , \forall x , 0<|x-x_0|<\delta \text{ 有}
f(x) \leq g(x) \leq h(x)\\
\text{且}\lim\limits_{x \to x_0}f(x) = \lim\limits_{x \to x_0}h(x) = A
\end{aligned}
$$
则
$$
\lim\limits_{x \to x_0}g(x) = A
$$

\end{document}

\end{document}
